% Options for packages loaded elsewhere
\PassOptionsToPackage{unicode}{hyperref}
\PassOptionsToPackage{hyphens}{url}
\documentclass[
]{article}
\usepackage{xcolor}
\usepackage{amsmath,amssymb}
\setcounter{secnumdepth}{-\maxdimen} % remove section numbering
\usepackage{iftex}
\ifPDFTeX
  \usepackage[T1]{fontenc}
  \usepackage[utf8]{inputenc}
  \usepackage{textcomp} % provide euro and other symbols
\else % if luatex or xetex
  \usepackage{unicode-math} % this also loads fontspec
  \defaultfontfeatures{Scale=MatchLowercase}
  \defaultfontfeatures[\rmfamily]{Ligatures=TeX,Scale=1}
\fi
\usepackage{lmodern}
\ifPDFTeX\else
  % xetex/luatex font selection
\fi
% Use upquote if available, for straight quotes in verbatim environments
\IfFileExists{upquote.sty}{\usepackage{upquote}}{}
\IfFileExists{microtype.sty}{% use microtype if available
  \usepackage[]{microtype}
  \UseMicrotypeSet[protrusion]{basicmath} % disable protrusion for tt fonts
}{}
\makeatletter
\@ifundefined{KOMAClassName}{% if non-KOMA class
  \IfFileExists{parskip.sty}{%
    \usepackage{parskip}
  }{% else
    \setlength{\parindent}{0pt}
    \setlength{\parskip}{6pt plus 2pt minus 1pt}}
}{% if KOMA class
  \KOMAoptions{parskip=half}}
\makeatother
\usepackage{color}
\usepackage{fancyvrb}
\newcommand{\VerbBar}{|}
\newcommand{\VERB}{\Verb[commandchars=\\\{\}]}
\DefineVerbatimEnvironment{Highlighting}{Verbatim}{commandchars=\\\{\}}
% Add ',fontsize=\small' for more characters per line
\newenvironment{Shaded}{}{}
\newcommand{\AlertTok}[1]{\textcolor[rgb]{1.00,0.00,0.00}{\textbf{#1}}}
\newcommand{\AnnotationTok}[1]{\textcolor[rgb]{0.38,0.63,0.69}{\textbf{\textit{#1}}}}
\newcommand{\AttributeTok}[1]{\textcolor[rgb]{0.49,0.56,0.16}{#1}}
\newcommand{\BaseNTok}[1]{\textcolor[rgb]{0.25,0.63,0.44}{#1}}
\newcommand{\BuiltInTok}[1]{\textcolor[rgb]{0.00,0.50,0.00}{#1}}
\newcommand{\CharTok}[1]{\textcolor[rgb]{0.25,0.44,0.63}{#1}}
\newcommand{\CommentTok}[1]{\textcolor[rgb]{0.38,0.63,0.69}{\textit{#1}}}
\newcommand{\CommentVarTok}[1]{\textcolor[rgb]{0.38,0.63,0.69}{\textbf{\textit{#1}}}}
\newcommand{\ConstantTok}[1]{\textcolor[rgb]{0.53,0.00,0.00}{#1}}
\newcommand{\ControlFlowTok}[1]{\textcolor[rgb]{0.00,0.44,0.13}{\textbf{#1}}}
\newcommand{\DataTypeTok}[1]{\textcolor[rgb]{0.56,0.13,0.00}{#1}}
\newcommand{\DecValTok}[1]{\textcolor[rgb]{0.25,0.63,0.44}{#1}}
\newcommand{\DocumentationTok}[1]{\textcolor[rgb]{0.73,0.13,0.13}{\textit{#1}}}
\newcommand{\ErrorTok}[1]{\textcolor[rgb]{1.00,0.00,0.00}{\textbf{#1}}}
\newcommand{\ExtensionTok}[1]{#1}
\newcommand{\FloatTok}[1]{\textcolor[rgb]{0.25,0.63,0.44}{#1}}
\newcommand{\FunctionTok}[1]{\textcolor[rgb]{0.02,0.16,0.49}{#1}}
\newcommand{\ImportTok}[1]{\textcolor[rgb]{0.00,0.50,0.00}{\textbf{#1}}}
\newcommand{\InformationTok}[1]{\textcolor[rgb]{0.38,0.63,0.69}{\textbf{\textit{#1}}}}
\newcommand{\KeywordTok}[1]{\textcolor[rgb]{0.00,0.44,0.13}{\textbf{#1}}}
\newcommand{\NormalTok}[1]{#1}
\newcommand{\OperatorTok}[1]{\textcolor[rgb]{0.40,0.40,0.40}{#1}}
\newcommand{\OtherTok}[1]{\textcolor[rgb]{0.00,0.44,0.13}{#1}}
\newcommand{\PreprocessorTok}[1]{\textcolor[rgb]{0.74,0.48,0.00}{#1}}
\newcommand{\RegionMarkerTok}[1]{#1}
\newcommand{\SpecialCharTok}[1]{\textcolor[rgb]{0.25,0.44,0.63}{#1}}
\newcommand{\SpecialStringTok}[1]{\textcolor[rgb]{0.73,0.40,0.53}{#1}}
\newcommand{\StringTok}[1]{\textcolor[rgb]{0.25,0.44,0.63}{#1}}
\newcommand{\VariableTok}[1]{\textcolor[rgb]{0.10,0.09,0.49}{#1}}
\newcommand{\VerbatimStringTok}[1]{\textcolor[rgb]{0.25,0.44,0.63}{#1}}
\newcommand{\WarningTok}[1]{\textcolor[rgb]{0.38,0.63,0.69}{\textbf{\textit{#1}}}}
\usepackage{longtable,booktabs,array}
\newcounter{none} % for unnumbered tables
\usepackage{calc} % for calculating minipage widths
% Correct order of tables after \paragraph or \subparagraph
\usepackage{etoolbox}
\makeatletter
\patchcmd\longtable{\par}{\if@noskipsec\mbox{}\fi\par}{}{}
\makeatother
% Allow footnotes in longtable head/foot
\IfFileExists{footnotehyper.sty}{\usepackage{footnotehyper}}{\usepackage{footnote}}
\makesavenoteenv{longtable}
\setlength{\emergencystretch}{3em} % prevent overfull lines
\providecommand{\tightlist}{%
  \setlength{\itemsep}{0pt}\setlength{\parskip}{0pt}}
\usepackage{bookmark}
\IfFileExists{xurl.sty}{\usepackage{xurl}}{} % add URL line breaks if available
\urlstyle{same}
% fallback for those not using the hyperref driver hyperxmp:
\makeatletter
\@ifundefined{xmpquote}{\newcommand{\xmpquote}[1]{#1}}{}
\makeatother
\hypersetup{
  hidelinks,
  pdfcreator={LaTeX via pandoc}}

\author{}
\date{}

\begin{document}

\section{Boolean 함수 공간의 S₄ 군론적 분류와 구조 인식 이론의
연결}\label{boolean-uxd568uxc218-uxacf5uxac04uxc758-sux2084-uxad70uxb860uxc801-uxbd84uxb958uxc640-uxad6cuxc870-uxc778uxc2dd-uxc774uxb860uxc758-uxc5f0uxacb0}

\begin{quote}
\textbf{관련 문서:} 향후 연구 전체 로드맵은
\href{TODO.md}{\texttt{TODO.md}}의 ``미래 연구(FutureWorks)'' 절 참조.\\
\textbf{수학 원본:} 상세 계산 과정 전체는
\href{theory/S4GroupAnalysis.md}{\texttt{theory/S4GroupAnalysis.md}}
참조.
\end{quote}

\textbf{저자:} 최종훈 (인하공업전문대학 출신 · 독립 연구자)\\
\textbf{작성일:} 2026-06-27\\
\textbf{저장소:}
\href{https://github.com/inha20/6BooleanFunctionSpaceTheory}{6BooleanFunctionSpaceTheory}\\
\textbf{연구 프로그램:} Structure Recognition Research Program (Branch
8)\\
\textbf{버전:} 1.0 (Phase 2 완료 기준)

\begin{center}\rule{0.5\linewidth}{0.5pt}\end{center}

\subsection{초록}\label{uxcd08uxb85d}

본 논문은 4변수 Boolean 함수 공간을 게이트 생성(Space(G)) 관점과 S₄
대칭군 관점에서 분류한 수학적 연구를 기술한다. 핵심 성과는 세 가지 주요
함수 공간에 대한 완전한 궤도 분류이다: Space(XOR)(아핀 함수 32개,
S₄-궤도 10개), Space(AND)(단조 함수 168개, S₄-궤도 30개),
Space(NOT)(자기쌍대 함수 256개, S₄-궤도 32개). 아울러 Space(NOT)에는
S₄-불동점이 존재하지 않음을 증명한다(정리 3.5.2). 이 수학적 결과들은
선행 연구 3편(Repo 1--3)의 경험적 발견을 군론 언어로 통일하며, 구조 인식
이론(SRT, Repo 4)의 핵심 가설(H3·H4·H8·H9)에 구체적인 수학적 토대를
제공한다. 분류의 정확성은 Burnside 보조정리에 의한 이론적 계산과 Python
컴퓨터 열거를 통한 독립 검증 양쪽으로 확인되었다.

\textbf{핵심어:} Boolean 함수 공간, S₄ 대칭군, Burnside 보조정리, 카르노
맵, Post 격자, 궤도 분류, 구조 인식 이론

\begin{center}\rule{0.5\linewidth}{0.5pt}\end{center}

\subsection{1. 서론}\label{uxc11cuxb860}

\subsubsection{1.1 연구의 배경}\label{uxc5f0uxad6cuxc758-uxbc30uxacbd}

\begin{quote}
\textbf{{[}연구자 보완 필요: 이 연구에 이르게 된 개인적 계기를 서술해
주세요.{]}}\\
\emph{예시 프롬프트: Karnaugh map을 처음 접한 시점, 체커보드 패턴에서
처음 ``왜?''라는 질문이 생긴 상황, 선행 연구들(Repo 1--3)을 진행하면서
군론적 통합의 필요성을 느끼게 된 계기 등}
\end{quote}

4변수 Boolean 함수 전체 공간 \(\mathcal{F}_4\)는
\(2^{16} = 65{,}536\)개의 함수를 포함한다. 이 방대한 공간 안에서
연구자가 ``주목할 만한 구조''를 발견하는 과정은 어떻게 일어나는가? 선행
연구들(Repo 1--3)은 이 질문에 대한 경험적 관찰을 축적해 왔다.

\begin{itemize}
\tightlist
\item
  \textbf{Repo 1} --- Karnaugh map에서 특정 함수들이 변수 배열에
  무관하게 동일한 패턴을 유지함을 발견
\item
  \textbf{Repo 2} --- 대칭 Boolean 함수들이 Hamming weight 레이어 구조를
  형성함을 관찰
\item
  \textbf{Repo 3} --- 변수 재배열(S₄ 치환) 하의 동치류를 체계적으로 분류
\end{itemize}

이 세 연구는 각각 독립적으로 수행되었으나, 하나의 공통 수학적
프레임워크로 통합될 수 있다는 예감이 이 연구의 출발점이다.

\subsubsection{1.2 연구 목표}\label{uxc5f0uxad6c-uxbaa9uxd45c}

본 연구는 두 가지 목표를 가진다.

\textbf{목표 1 (수학적).} 4변수 Boolean 함수 공간을 게이트 생성 관점과
S₄ 대칭군 관점에서 완전히 분류한다. 구체적으로, 주요 게이트(XOR, AND,
NOT)로 생성되는 함수 공간 Space(G) 각각에 대해 S₄ 작용에 의한 궤도 수와
각 궤도의 구조를 결정한다.

\textbf{목표 2 (이론적).} 분류 결과를 구조 인식 이론(SRT, Repo 4)의 핵심
가설들과 연결한다. 특히, 수학적 분류가 ``인간이 어떤 함수에 먼저 주의를
기울이는가''(H3)와 ``어떤 함수를 연구 가치 있다고 판단하는가''(H4)를
설명하는 데 어떻게 기여하는지를 논의한다.

\subsubsection{1.3 논문의 구성}\label{uxb17cuxbb38uxc758-uxad6cuxc131}

2장은 이론적 배경(Space(G) 정의, Post 격자, 카르노 맵 기하학)을
정리한다. 3장은 핵심 수학적 성과인 S₄ 궤도 분류를 공간별로 전개한다.
4장은 선행 연구(Repo 1--3)와의 통합 관계를 정식화한다. 5장은 SRT
가설과의 연결을 논의한다. 6장은 이 연구에 이르기까지의 연구 기원
타임라인을 정리한다.

\begin{center}\rule{0.5\linewidth}{0.5pt}\end{center}

\subsection{2. 이론적 배경}\label{uxc774uxb860uxc801-uxbc30uxacbd}

\subsubsection{2.1 4변수 Boolean 함수
공간}\label{uxbcc0uxc218-boolean-uxd568uxc218-uxacf5uxac04}

\textbf{정의 2.1 (전체 함수 공간).}

\[\mathcal{F}_4 = \{ f : \{0,1\}^4 \to \{0,1\} \}\]

이 공간의 크기는 \(|\mathcal{F}_4| = 2^{2^4} = 2^{16} = 65{,}536\)이다.
각 함수는 16비트 진리표 벡터, 4×4 카르노 맵, 또는 게이트 회로로
표현된다.

\subsubsection{2.2 게이트 생성 함수 공간
Space(G)}\label{uxac8cuxc774uxd2b8-uxc0dduxc131-uxd568uxc218-uxacf5uxac04-spaceg}

단일 게이트 G에 대해, G만을 사용하여 구성 가능한 모든 Boolean 함수의
집합을 \textbf{Space(G)}로 정의한다.

\textbf{정의 2.2 (Space(G)).}

\[\text{Space}(G) = \{ f \in \mathcal{F}_4 \mid f\text{는 게이트 }G\text{와 입력 변수, 상수 }\{0,1\}\text{만으로 구성 가능} \}\]

이는 대수학의 \textbf{클론(clone)} 개념과 동치이다:
\(\text{Space}(G) = [G]\) (G가 생성하는 Post 클론). 주요 게이트별
Space(G)의 요약은 다음과 같다.

{\def\LTcaptype{none} % do not increment counter
\begin{longtable}[]{@{}llll@{}}
\toprule\noalign{}
게이트 G & Space(G) & 크기 (4변수) & Post 격자 위치 \\
\midrule\noalign{}
\endhead
\bottomrule\noalign{}
\endlastfoot
NAND & 전체 Boolean 공간 (Ω) & 65,536 & 최상위 \\
NOR & 전체 Boolean 공간 (Ω) & 65,536 & 최상위 \\
XOR & 아핀(Affine) 함수 공간 (L) & 32 & 선형 클론 \\
AND & 단조(Monotone) 함수 공간 (M) & 168 & 단조 클론 \\
NOT & 자기쌍대(Self-dual) 함수 공간 (D) & 256 & 쌍대 클론 \\
\end{longtable}
}

\textbf{명제 2.1 (Space(G)의 S₄-불변성).} 임의의 게이트 G에 대해,
Space(G)는 S₄ 작용 하에 닫혀 있다. 즉, \(f \in \text{Space}(G)\)이고
\(\sigma \in S_4\)이면 \(f \circ \sigma \in \text{Space}(G)\).

\emph{증명 스케치.} f가 G-게이트만으로 구현되는 회로를 가진다면,
\(f \circ \sigma\)는 같은 회로에서 입력 와이어를 σ에 따라 재배치한
것이다. 회로의 게이트 종류(G)는 변하지 않으므로
\(f \circ \sigma \in \text{Space}(G)\). ∎

이 명제의 결론으로, \textbf{각 Space(G)는 S₄ 궤도들의 합집합으로
분해}된다.

\subsubsection{2.3 Post의 격자와
만능성}\label{postuxc758-uxaca9uxc790uxc640-uxb9ccuxb2a5uxc131}

\textbf{정리 2.1 (Post, 1941).} Boolean 함수 집합 \(F\)가 함수적으로
완전하기 위한 필요충분조건은 F가 다음 5가지 Post 클래스 중 어느 것에도
포함되지 않는 것이다.

{\def\LTcaptype{none} % do not increment counter
\begin{longtable}[]{@{}lll@{}}
\toprule\noalign{}
클래스 & 조건 & Space(G) 대응 \\
\midrule\noalign{}
\endhead
\bottomrule\noalign{}
\endlastfoot
T₀ & 0을 보존 & Space(AND) ⊄ T₀ (상수 1 포함) \\
T₁ & 1을 보존 & Space(AND) ⊄ T₁ (상수 0 포함) \\
S & 자기쌍대 & Space(NOT) = D \\
M & 단조 & Space(AND) = M (상수 0, 1 포함, 168개) \\
L & 선형(아핀) & Space(XOR) = L \\
\end{longtable}
}

NAND는 5가지 클래스 중 어느 것도 보존하지 않으므로 만능 게이트이다.

\subsubsection{2.4 카르노 맵의 기하학적
해석}\label{uxce74uxb974uxb178-uxb9f5uxc758-uxae30uxd558uxd559uxc801-uxd574uxc11d}

4변수 카르노 맵은 단순한 회로 최소화 도구가 아니라, \(\mathcal{F}_4\)를
기하학적으로 시각화하는 공간이다.

\begin{verbatim}
        CD
        00  01  11  10
AB  00 | 0 | 1 | 3 | 2 |
    01 | 4 | 5 | 7 | 6 |
    11 |12 |13 |15 |14 |
    10 | 8 | 9 |11 |10 |
\end{verbatim}

\begin{itemize}
\tightlist
\item
  인접 칸은 \textbf{정확히 한 변수만 다름} (Gray 코드 순서)
\item
  맵은 \textbf{환면체(Torus)} 위에 놓인 것으로 볼 수 있음: 상단 행과
  하단 행(AB=00과 AB=10)이 연결되고, 좌측 열과 우측 열(CD=00과 CD=10)이
  연결된다
\item
  따라서 ``모서리'', ``좌상단'', ``우하단'' 등 \textbf{절대 위치 표현은
  토러스 위상 하에서 의미를 잃는다} --- 어떤 셀도 다른 셀보다 특별히
  ``끝''에 있지 않다
\item
  65,536개 Boolean 함수 = 4×4 격자의 서로 다른 0/1 채색
\end{itemize}

\textbf{토러스 구조의 주요 함의:}

{\def\LTcaptype{none} % do not increment counter
\begin{longtable}[]{@{}
  >{\raggedright\arraybackslash}p{(\linewidth - 4\tabcolsep) * \real{0.4722}}
  >{\raggedright\arraybackslash}p{(\linewidth - 4\tabcolsep) * \real{0.3889}}
  >{\raggedright\arraybackslash}p{(\linewidth - 4\tabcolsep) * \real{0.1389}}@{}}
\toprule\noalign{}
\begin{minipage}[b]{\linewidth}\raggedright
카르노 맵상 인접 쌍
\end{minipage} & \begin{minipage}[b]{\linewidth}\raggedright
토러스 연결 여부
\end{minipage} & \begin{minipage}[b]{\linewidth}\raggedright
설명
\end{minipage} \\
\midrule\noalign{}
\endhead
\bottomrule\noalign{}
\endlastfoot
AB=00 ↔ AB=10 (상하 끝 행) & ✅ 연결됨 & Gray 코드: 00과 10은 한 비트
차이 \\
CD=00 ↔ CD=10 (좌우 끝 열) & ✅ 연결됨 & Gray 코드: 00과 10은 한 비트
차이 \\
네 모서리 (m0, m2, m8, m10) & 모두 인접 & 토러스에서 한 변수씩만 다른
이웃을 가짐 \\
\end{longtable}
}

따라서 같은 궤도에 속하는 함수들은 카르노 맵에서 \textbf{회전·반사
대칭인 패턴}을 가지며, 특정 시각적 패턴이 맵의 어느 위치에 나타나느냐는
변수 명명 방식(배열 순서)에만 의존한다.

Space(G)에 속하는 함수들은 카르노 맵 위에서 공통 시각적 시그니처를
보인다. 아핀 함수는 체커보드·줄무늬 패턴을, 단조 함수는 ``위쪽이
닫힌(upward-closed)'' 패턴을(단, Gray 코드 배열로 인해 시각적으로
직관적이지 않을 수 있음), 자기쌍대 함수는 180° 회전 시 0↔1이 교환되는
패턴을 형성한다.

\subsubsection{2.5 S₄의 Boolean 함수에 대한
작용}\label{sux2084uxc758-boolean-uxd568uxc218uxc5d0-uxb300uxd55c-uxc791uxc6a9}

\textbf{군:} \(S_4\) = 4개 원소 \(\{A, B, C, D\}\)의 모든 순열의 군
(크기 24)

\textbf{작용:}

\[(\sigma \cdot f)(A, B, C, D) = f(\sigma^{-1}(A), \sigma^{-1}(B), \sigma^{-1}(C), \sigma^{-1}(D))\]

S₄의 5가지 켤레류(conjugacy class):

{\def\LTcaptype{none} % do not increment counter
\begin{longtable}[]{@{}lll@{}}
\toprule\noalign{}
켤레류 & 원소 수 & 예시 \\
\midrule\noalign{}
\endhead
\bottomrule\noalign{}
\endlastfoot
항등원 & 1 & (A)(B)(C)(D) \\
호환 (transposition) & 6 & (AB), (AC), \ldots{} \\
3-순환 (3-cycle) & 8 & (ABC), (ABD), \ldots{} \\
4-순환 (4-cycle) & 6 & (ABCD), (ABDC), \ldots{} \\
이중 호환 & 3 & (AB)(CD), (AC)(BD), (AD)(BC) \\
\textbf{합계} & \textbf{24} & \\
\end{longtable}
}

\textbf{Burnside의 보조정리.} S₄-궤도 수를 계산하는 핵심 도구:

\[|\text{Orbits}| = \frac{1}{|S_4|} \sum_{\sigma \in S_4} |\text{Fix}(\sigma, \text{Space}(G))|\]

여기서
\(\text{Fix}(\sigma, \text{Space}(G)) = \{ f \in \text{Space}(G) : f \circ \sigma = f \}\)이다.

\begin{center}\rule{0.5\linewidth}{0.5pt}\end{center}

\subsection{3. 핵심 수학적 성과: S₄ 궤도 완전
분류}\label{uxd575uxc2ec-uxc218uxd559uxc801-uxc131uxacfc-sux2084-uxada4uxb3c4-uxc644uxc804-uxbd84uxb958}

\subsubsection{3.1 Space(XOR) = L (아핀 클론): 10개
궤도}\label{spacexor-l-uxc544uxd540-uxd074uxb860-10uxac1c-uxada4uxb3c4}

\paragraph{3.1.1 구조}\label{uxad6cuxc870}

4변수 아핀 함수:

\[f(A, B, C, D) = a_0 \oplus a_1 A \oplus a_2 B \oplus a_3 C \oplus a_4 D, \quad a_i \in \{0, 1\}\]

총 원소 수: \(2^5 = 32\)개.

\textbf{S₄ 작용의 효과:} 순열 σ가 변수 \(x_i\)와 \(x_j\)를 교환하면,
불동점 조건은 \(a_i = a_j\).

\paragraph{3.1.2 Burnside 계산}\label{burnside-uxacc4uxc0b0}

{\def\LTcaptype{none} % do not increment counter
\begin{longtable}[]{@{}llll@{}}
\toprule\noalign{}
켤레류 & 원소 수 & σ의 궤도 구조 & Fix(σ, L) \\
\midrule\noalign{}
\endhead
\bottomrule\noalign{}
\endlastfoot
항등원 & 1 & 4개 1-궤도 & \(2^5 = 32\) \\
호환 (AB)(C)(D) & 6 & \{A,B\} + \{C\} + \{D\} & \(2^4 = 16\) \\
3-순환 (ABC)(D) & 8 & \{A,B,C\} + \{D\} & \(2^3 = 8\) \\
4-순환 (ABCD) & 6 & \{A,B,C,D\} & \(2^2 = 4\) \\
이중 호환 (AB)(CD) & 3 & \{A,B\} + \{C,D\} & \(2^3 = 8\) \\
\end{longtable}
}

\[|\text{Orbits}(S_4, L)| = \frac{1}{24}(32 + 6 \times 16 + 8 \times 8 + 6 \times 4 + 3 \times 8) = \frac{240}{24} = \mathbf{10}\]

\paragraph{3.1.3 10개 궤도 완전
열거}\label{uxac1c-uxada4uxb3c4-uxc644uxc804-uxc5f4uxac70}

{\def\LTcaptype{none} % do not increment counter
\begin{longtable}[]{@{}llll@{}}
\toprule\noalign{}
궤도 & 대표 함수 & 크기 & 의미 \\
\midrule\noalign{}
\endhead
\bottomrule\noalign{}
\endlastfoot
O₁ & \(0\) & 1 & 상수 0 \\
O₂ & \(1\) & 1 & 상수 1 \\
O₃ & \(A\) & 4 & 단일 변수 투영 \\
O₄ & \(A \oplus 1\) & 4 & 단일 변수 투영의 보수 \\
O₅ & \(A \oplus B\) & 6 & 2변수 XOR \\
O₆ & \(A \oplus B \oplus 1\) & 6 & 2변수 XOR의 보수 \\
O₇ & \(A \oplus B \oplus C\) & 4 & 3변수 XOR \\
O₈ & \(A \oplus B \oplus C \oplus 1\) & 4 & 3변수 XOR의 보수 \\
O₉ & \(A \oplus B \oplus C \oplus D\) & 1 & \textbf{완전 패리티
(체커보드)} \\
O₁₀ & \(A \oplus B \oplus C \oplus D \oplus 1\) & 1 & 체커보드의 보수 \\
\end{longtable}
}

\textbf{검증:} \(1+1+4+4+6+6+4+4+1+1 = 32\) ✓

\paragraph{3.1.4 S₄-불동점}\label{sux2084-uxbd88uxb3d9uxc810}

S₄의 모든 원소에 의해 고정되는 함수(4-순환 조건:
\(a_1 = a_2 = a_3 = a_4\)):

{\def\LTcaptype{none} % do not increment counter
\begin{longtable}[]{@{}
  >{\raggedright\arraybackslash}p{(\linewidth - 6\tabcolsep) * \real{0.1622}}
  >{\raggedright\arraybackslash}p{(\linewidth - 6\tabcolsep) * \real{0.5676}}
  >{\raggedright\arraybackslash}p{(\linewidth - 6\tabcolsep) * \real{0.1351}}
  >{\raggedright\arraybackslash}p{(\linewidth - 6\tabcolsep) * \real{0.1351}}@{}}
\toprule\noalign{}
\begin{minipage}[b]{\linewidth}\raggedright
\(a_0\)
\end{minipage} & \begin{minipage}[b]{\linewidth}\raggedright
\(a_1(=a_2=a_3=a_4)\)
\end{minipage} & \begin{minipage}[b]{\linewidth}\raggedright
함수
\end{minipage} & \begin{minipage}[b]{\linewidth}\raggedright
의미
\end{minipage} \\
\midrule\noalign{}
\endhead
\bottomrule\noalign{}
\endlastfoot
0 & 0 & \(f = 0\) & 상수 0 \\
1 & 0 & \(f = 1\) & 상수 1 \\
0 & 1 & \(A \oplus B \oplus C \oplus D\) & \textbf{체커보드 (Repo 1의
최초 발견)} \\
1 & 1 & \(A \oplus B \oplus C \oplus D \oplus 1\) & 체커보드 보수 \\
\end{longtable}
}

\textbf{결론:} Space(XOR)의 S₄-불동점은 정확히 \textbf{4개}이며, 이 중
체커보드 패턴(\(A \oplus B \oplus C \oplus D\))은 Repo 1의 연구 출발점이
된 함수이다.

\begin{center}\rule{0.5\linewidth}{0.5pt}\end{center}

\subsubsection{3.2 Space(AND) = M (단조 클론): 30개
궤도}\label{spaceand-m-uxb2e8uxc870-uxd074uxb860-30uxac1c-uxada4uxb3c4}

\paragraph{3.2.1 구조}\label{uxad6cuxc870-1}

4변수 단조 함수 수: \textbf{Dedekind 수 D(4) = 168}.

단조 함수는 \textbf{반사슬(antichain)}
\(\mathcal{A} \subseteq 2^{\{A,B,C,D\}}\)로 유일하게 특성화된다:
\(f_{\mathcal{A}}(x) = 1\) iff \(\exists S \in \mathcal{A}\) s.t.
\(S \subseteq \{i : x_i = 1\}\).

\paragraph{3.2.2 Burnside 계산}\label{burnside-uxacc4uxc0b0-1}

{\def\LTcaptype{none} % do not increment counter
\begin{longtable}[]{@{}lll@{}}
\toprule\noalign{}
켤레류 & 원소 수 & Fix(σ, M) \\
\midrule\noalign{}
\endhead
\bottomrule\noalign{}
\endlastfoot
항등원 & 1 & 168 \\
호환 (AB) & 6 & \textbf{50} \\
3-순환 (ABC) & 8 & \textbf{15} \\
4-순환 (ABCD) & 6 & \textbf{8} \\
이중 호환 (AB)(CD) & 3 & \textbf{28} \\
\end{longtable}
}

\[|\text{Orbits}(S_4, M)| = \frac{1}{24}(168 + 6 \times 50 + 8 \times 15 + 6 \times 8 + 3 \times 28) = \frac{720}{24} = \mathbf{30}\]

검증: Python 직접 열거(컴퓨터 계산)로 30개 독립 확인, 크기 합 = 168 ✓

\paragraph{3.2.3 S₄-불동점: 임계값 함수
6개}\label{sux2084-uxbd88uxb3d9uxc810-uxc784uxacc4uxac12-uxd568uxc218-6uxac1c}

S₄ 전체에 의해 고정되는 단조 함수는 입력의 \textbf{Hamming
가중치(HW)}에만 의존하는 함수, 즉 임계값 함수(threshold function)이다.

\[\text{Fix}(S_4, \text{Space(AND)}) = \{0,\ A \vee B \vee C \vee D,\ \text{Maj}_2,\ \text{Maj}_3,\ A \wedge B \wedge C \wedge D,\ 1\}\]

이 6개 함수는 Repo 2(대칭 Boolean 함수)가 다루는 대칭∧단조 함수의 정확한
목록이다.

\paragraph{3.2.4 궤도 크기
분포}\label{uxada4uxb3c4-uxd06cuxae30-uxbd84uxd3ec}

{\def\LTcaptype{none} % do not increment counter
\begin{longtable}[]{@{}lll@{}}
\toprule\noalign{}
궤도 크기 & 궤도 수 & 소계 \\
\midrule\noalign{}
\endhead
\bottomrule\noalign{}
\endlastfoot
1 & 6 & 6 \\
3 & 2 & 6 \\
4 & 9 & 36 \\
6 & 6 & 36 \\
12 & 7 & 84 \\
\textbf{합계} & \textbf{30} & \textbf{168 ✓} \\
\end{longtable}
}

\paragraph{3.2.5 30개 궤도 완전 목록 (S4GroupAnalysis.md OC-2
발췌)}\label{uxac1c-uxada4uxb3c4-uxc644uxc804-uxbaa9uxb85d-s4groupanalysis.md-oc-2-uxbc1cuxcdcc}

\begin{quote}
\emph{완전한 목록은
\href{theory/S4GroupAnalysis.md}{\texttt{theory/S4GroupAnalysis.md}} §5
OC-2 참조.}
\end{quote}

{\def\LTcaptype{none} % do not increment counter
\begin{longtable}[]{@{}
  >{\raggedright\arraybackslash}p{(\linewidth - 8\tabcolsep) * \real{0.1471}}
  >{\raggedright\arraybackslash}p{(\linewidth - 8\tabcolsep) * \real{0.1765}}
  >{\raggedright\arraybackslash}p{(\linewidth - 8\tabcolsep) * \real{0.2647}}
  >{\raggedright\arraybackslash}p{(\linewidth - 8\tabcolsep) * \real{0.2647}}
  >{\raggedright\arraybackslash}p{(\linewidth - 8\tabcolsep) * \real{0.1471}}@{}}
\toprule\noalign{}
\begin{minipage}[b]{\linewidth}\raggedright
궤도
\end{minipage} & \begin{minipage}[b]{\linewidth}\raggedright
크기
\end{minipage} & \begin{minipage}[b]{\linewidth}\raggedright
대표 함수
\end{minipage} & \begin{minipage}[b]{\linewidth}\raggedright
S₄-불동점
\end{minipage} & \begin{minipage}[b]{\linewidth}\raggedright
비고
\end{minipage} \\
\midrule\noalign{}
\endhead
\bottomrule\noalign{}
\endlastfoot
\textbf{O01} & 1 & \(0\) & ✓ & 임계값 함수 \\
\textbf{O02} & 1 & \(A \wedge B \wedge C \wedge D\) & ✓ & Threshold-4 \\
\textbf{O03} & 1 & \(\text{Maj}_3\) (at-least-3) & ✓ & Threshold-3 \\
\textbf{O04} & 1 & \(\text{Maj}_2\) (at-least-2) & ✓ & Threshold-2 \\
\textbf{O05} & 1 & \(A \vee B \vee C \vee D\) & ✓ & Threshold-1 \\
\textbf{O06} & 1 & \(1\) & ✓ & 임계값 함수 \\
\textbf{O07} & 3 & \((A \wedge B) \vee (C \wedge D)\) & ✗ & 쌍 AND-OR
(3가지 변수 쌍) \\
\textbf{O08} & 3 & \((A \vee B) \wedge (C \vee D)\) & ✗ & 쌍 OR-AND
(3가지 변수 쌍) \\
\textbf{O09} & 4 & \(B \wedge C \wedge D\) & ✗ & 3-AND (4가지 변수
선택) \\
\textbf{O10} & 4 & \(B \wedge C \wedge D \vee A \wedge C \wedge D \vee A \wedge B \wedge D\) & ✗ & 3개 3-AND의 합 \\
\textbf{O11} & 4 & \(C \wedge D \vee B \wedge D \vee A \wedge D\) & ✗ &
한 변수 고정 OR \\
\textbf{O12} & 4 & \(D\) & ✗ & 단일 변수 투영 \\
\textbf{O13} & 4 & \(C \wedge D \vee B \wedge D \vee B \wedge C\) & ✗
& \\
\textbf{O14} & 4 &
\(C \wedge D \vee B \wedge D \vee A \wedge D \vee A \wedge B \wedge C\)
& ✗ & \\
\textbf{O15} & 4 & \(D \vee A \wedge B \wedge C\) & ✗ & \\
\textbf{O16} & 4 & \(D \vee B \wedge C \vee A \wedge C \vee A \wedge B\)
& ✗ & \\
\textbf{O17} & 4 & \(D \vee C \vee B\) & ✗ & 3-OR \\
\textbf{O18} & 6 & \(B \wedge C \wedge D \vee A \wedge C \wedge D\) & ✗
& \\
\textbf{O19} & 6 & \(C \wedge D\) & ✗ & 2-AND \\
\textbf{O20} & 6 &
\(C \wedge D \vee A \wedge B \wedge D \vee A \wedge B \wedge C\) & ✗
& \\
\textbf{O21} & 6 &
\(C \wedge D \vee B \wedge D \vee B \wedge C \vee A \wedge D \vee A \wedge C\)
& ✗ & \\
\textbf{O22} & 6 & \(D \vee C\) & ✗ & 2-OR \\
\textbf{O23} & 6 & \(D \vee C \vee A \wedge B\) & ✗ & \\
\textbf{O24} & 12 & \(C \wedge D \vee A \wedge B \wedge D\) & ✗ & \\
\textbf{O25} & 12 & \(C \wedge D \vee B \wedge D\) & ✗ &
\(A \wedge (B \vee C)\) 류 \\
\textbf{O26} & 12 &
\(C \wedge D \vee B \wedge D \vee A \wedge B \wedge C\) & ✗ & \\
\textbf{O27} & 12 & \(C \wedge D \vee B \wedge C \vee A \wedge D\) & ✗
& \\
\textbf{O28} & 12 &
\(C \wedge D \vee B \wedge D \vee B \wedge C \vee A \wedge D\) & ✗ & \\
\textbf{O29} & 12 & \(D \vee B \wedge C\) & ✗ & \((A \wedge B) \vee C\)
류 \\
\textbf{O30} & 12 & \(D \vee B \wedge C \vee A \wedge C\) & ✗ & \\
\end{longtable}
}

\textbf{크기 검증:}
\(6 \times 1 + 2 \times 3 + 9 \times 4 + 6 \times 6 + 7 \times 12 = 6+6+36+36+84 = \mathbf{168}\)
✓

\textbf{SRT 연결:} S₄-불동점(크기 1, O01\textasciitilde O06)이 Repo
1·2에서 처음 관찰된 함수들과 정확히 일치한다. 궤도 크기가 커질수록(O07
크기 3 → O30 크기 12) 대칭성이 낮아지며, H3 예측에 따르면 주의를 끄는
빈도도 낮아진다.

\begin{center}\rule{0.5\linewidth}{0.5pt}\end{center}

\subsubsection{3.3 Space(NOT) = D (자기쌍대 클론): 32개 궤도 및 불동점
부재}\label{spacenot-d-uxc790uxae30uxc30duxb300-uxd074uxb860-32uxac1c-uxada4uxb3c4-uxbc0f-uxbd88uxb3d9uxc810-uxbd80uxc7ac}

\paragraph{3.3.1 구조}\label{uxad6cuxc870-2}

\textbf{자기쌍대 함수:} \(f(x) = \neg f(\neg x)\) for all
\(x \in \{0,1\}^4\).

4변수에서 입력 16개는 8쌍의 \textbf{보수 쌍(complementary pairs)}
\(\{n, \overline{n}\}\) (\(\overline{n} = 15 - n\))으로 분류된다:

\[\{0,15\},\ \{1,14\},\ \{2,13\},\ \{3,12\},\ \{4,11\},\ \{5,10\},\ \{6,9\},\ \{7,8\}\]

각 쌍에서 정확히 하나에 출력 1을 배정해야 하므로, 자기쌍대 함수는
\(2^8 = \mathbf{256}\)개이다.

\textbf{명제 3.1.} Space(NOT)는 S₄ 작용 하에 닫혀 있다. ∎

\paragraph{3.3.2 핵심 관찰: 충돌
메커니즘}\label{uxd575uxc2ec-uxad00uxcc30-uxcda9uxb3cc-uxba54uxcee4uxb2c8uxc998}

σ의 입력 궤도와 보수 쌍 \(\{n, \overline{n}\}\)이 \textbf{같은 궤도}를
이루면, σ-대칭 조건 \(f(n) = f(\overline{n})\)과 자기쌍대 조건
\(f(n) = \neg f(\overline{n})\)이 동시에 요구되어 \textbf{해가 없다}
(Fix = 0).

\paragraph{3.3.3 Burnside 계산}\label{burnside-uxacc4uxc0b0-2}

{\def\LTcaptype{none} % do not increment counter
\begin{longtable}[]{@{}llll@{}}
\toprule\noalign{}
켤레류 & 원소 수 & Fix(σ, D) & 충돌 여부 \\
\midrule\noalign{}
\endhead
\bottomrule\noalign{}
\endlastfoot
항등원 & 1 & \textbf{256} & 없음 \\
호환 (AB) & 6 & \textbf{64} & 없음 \\
3-순환 (ABC) & 8 & \textbf{16} & 없음 \\
4-순환 (ABCD) & 6 & \textbf{0} & ✗ \(\{5,10\}\) 충돌 \\
이중 호환 (AB)(CD) & 3 & \textbf{0} & ✗ \(\{5,10\}, \{6,9\}\) 충돌 \\
\end{longtable}
}

\[|\text{Orbits}(S_4, D)| = \frac{1}{24}(256 + 6 \times 64 + 8 \times 16 + 0 + 0) = \frac{768}{24} = \mathbf{32}\]

검증: Python 직접 열거로 32개 독립 확인, 크기 합 = 256 ✓

\paragraph{3.3.4 정리: S₄-불동점의
부재}\label{uxc815uxb9ac-sux2084-uxbd88uxb3d9uxc810uxc758-uxbd80uxc7ac}

\textbf{정리 3.1 (Space(NOT)의 S₄-불동점 = 0).}\\
Space(NOT)에는 S₄-불동점이 존재하지 않는다.

\emph{증명.} S₄-불동점은 Hamming 가중치에만 의존해야 한다. HW=2
레이어에서 \(\{5, 10\}\) (5=0101, 10=1010, HW=2)은 보수 쌍이다. Hamming
가중치 의존성이면 \(f(5) = f(10)\)이지만, 자기쌍대 조건은
\(f(5) = \neg f(10)\)을 요구한다. 모순. ∎

\textbf{함의:} Space(AND)의 6개 불동점(임계값 함수)은 Space(NOT) 안에
존재하지 않는다. ``자기쌍대이면서 완전히 대칭인'' 함수는 4변수에서
존재할 수 없다.

\paragraph{3.3.5 궤도 크기
분포}\label{uxada4uxb3c4-uxd06cuxae30-uxbd84uxd3ec-1}

{\def\LTcaptype{none} % do not increment counter
\begin{longtable}[]{@{}lll@{}}
\toprule\noalign{}
궤도 크기 & 궤도 수 & 소계 \\
\midrule\noalign{}
\endhead
\bottomrule\noalign{}
\endlastfoot
4 & 16 & 64 \\
12 & 16 & 192 \\
\textbf{합계} & \textbf{32} & \textbf{256 ✓} \\
\end{longtable}
}

\begin{center}\rule{0.5\linewidth}{0.5pt}\end{center}

\subsubsection{3.4 L ∩ M (아핀이면서 단조): 3개
궤도}\label{l-m-uxc544uxd540uxc774uxba74uxc11c-uxb2e8uxc870-3uxac1c-uxada4uxb3c4}

아핀 함수 \(f = a_0 \oplus \bigoplus_i a_i x_i\)가 단조이기 위한 조건을
만족하는 함수:

\[L \cap M = \{0,\ 1,\ A,\ B,\ C,\ D\}\]

총 \textbf{6개} (Python 열거로 확인 ✓).

\textbf{비포함 예시:} \(A \oplus B\)는 아핀이지만
비단조(\(A=1,B=0 \to 1\), \(A=1,B=1 \to 0\)).

\paragraph{3.4.1 S₄-궤도 (3개)}\label{sux2084-uxada4uxb3c4-3uxac1c}

{\def\LTcaptype{none} % do not increment counter
\begin{longtable}[]{@{}llll@{}}
\toprule\noalign{}
궤도 & 원소 & 크기 & S₄-불동점 \\
\midrule\noalign{}
\endhead
\bottomrule\noalign{}
\endlastfoot
O₁ & \(\{0\}\) & 1 & ✓ \\
O₂ & \(\{1\}\) & 1 & ✓ \\
O₃ & \(\{A, B, C, D\}\) & 4 & ✗ \\
\end{longtable}
}

\[|\text{Orbits}(S_4, L \cap M)| = \frac{72}{24} = \mathbf{3} \quad \checkmark\]

\textbf{결론:} 단조이면서 아핀이고 S₄-불동점인 함수 = 상수함수 2개 뿐.

\begin{center}\rule{0.5\linewidth}{0.5pt}\end{center}

\subsubsection{3.5 전체 계층 구조
요약}\label{uxc804uxccb4-uxacc4uxce35-uxad6cuxc870-uxc694uxc57d}

\begin{verbatim}
전체 65,536개 함수
└── Space(NAND) = Ω (65,536개)
    ├── Space(XOR) = L (32개, 10궤도)
    │   ├── O₁ = {0} (1개)          ← S₄-불동점
    │   ├── O₂ = {1} (1개)          ← S₄-불동점
    │   ├── O₃ = {A,B,C,D} (4개)
    │   ├── O₄ = {A⊕1,...} (4개)
    │   ├── O₅ = {A⊕B,...} (6개)
    │   ├── O₆ = {A⊕B⊕1,...} (6개)
    │   ├── O₇ = {A⊕B⊕C,...} (4개)
    │   ├── O₈ = {A⊕B⊕C⊕1,...} (4개)
    │   ├── O₉ = {A⊕B⊕C⊕D} (1개)  ← S₄-불동점 · Repo 1 출발점
    │   └── O₁₀ = {¬(A⊕B⊕C⊕D)} (1개) ← S₄-불동점
    │
    ├── Space(AND) = M (168개, 30궤도)
    │   ├── O01 = {0} (크기 1)       ← S₄-불동점
    │   ├── O02 = {A∧B∧C∧D} (크기 1) ← S₄-불동점
    │   ├── O03 = {Maj₃} (크기 1)   ← S₄-불동점
    │   ├── O04 = {Maj₂} (크기 1)   ← S₄-불동점
    │   ├── O05 = {A∨B∨C∨D} (크기 1) ← S₄-불동점
    │   ├── O06 = {1} (크기 1)       ← S₄-불동점
    │   └── O07~O30: 나머지 162개 함수 (크기 3~12)
    │
    ├── Space(NOT) = D (256개, 32궤도)
    │   ├── 크기 4 궤도 16개 (64개 함수)
    │   └── 크기 12 궤도 16개 (192개 함수)
    │   ※ S₄-불동점 없음 (정리 3.1)
    │
    └── L ∩ M = {0, 1, A, B, C, D} (6개, 3궤도)
\end{verbatim}

\textbf{분류 결과 요약:}

{\def\LTcaptype{none} % do not increment counter
\begin{longtable}[]{@{}lllll@{}}
\toprule\noalign{}
공간 & 함수 수 & S₄-궤도 수 & S₄-불동점 수 & 검증 방법 \\
\midrule\noalign{}
\endhead
\bottomrule\noalign{}
\endlastfoot
Space(XOR) = L & 32 & \textbf{10} & 4 & Burnside + Python ✓ \\
Space(AND) = M & 168 & \textbf{30} & 6 & Burnside + Python ✓ \\
Space(NOT) = D & 256 & \textbf{32} & \textbf{0} & Burnside + Python ✓ \\
L ∩ M & 6 & \textbf{3} & 2 & 직접 열거 + Burnside ✓ \\
\end{longtable}
}

\subsubsection{3.6 Python 컴퓨터 계산 검증
개요}\label{python-uxcef4uxd4e8uxd130-uxacc4uxc0b0-uxac80uxc99d-uxac1cuxc694}

\begin{Shaded}
\begin{Highlighting}[]
\ImportTok{from}\NormalTok{ itertools }\ImportTok{import}\NormalTok{ permutations}

\KeywordTok{def}\NormalTok{ apply\_permutation(f\_bits, perm):}
    \CommentTok{"""S4 원소를 Boolean 함수 (16비트 튜플)에 적용"""}
\NormalTok{    n\_vars }\OperatorTok{=} \DecValTok{4}
\NormalTok{    result }\OperatorTok{=}\NormalTok{ [}\DecValTok{0}\NormalTok{] }\OperatorTok{*}\NormalTok{ (}\DecValTok{2}\OperatorTok{**}\NormalTok{n\_vars)}
    \ControlFlowTok{for}\NormalTok{ minterm }\KeywordTok{in} \BuiltInTok{range}\NormalTok{(}\DecValTok{2}\OperatorTok{**}\NormalTok{n\_vars):}
\NormalTok{        bits }\OperatorTok{=}\NormalTok{ [(minterm }\OperatorTok{\textgreater{}\textgreater{}}\NormalTok{ (n\_vars}\OperatorTok{{-}}\DecValTok{1}\OperatorTok{{-}}\NormalTok{i)) }\OperatorTok{\&} \DecValTok{1} \ControlFlowTok{for}\NormalTok{ i }\KeywordTok{in} \BuiltInTok{range}\NormalTok{(n\_vars)]}
\NormalTok{        new\_bits }\OperatorTok{=}\NormalTok{ [bits[perm[i]] }\ControlFlowTok{for}\NormalTok{ i }\KeywordTok{in} \BuiltInTok{range}\NormalTok{(n\_vars)]}
\NormalTok{        new\_minterm }\OperatorTok{=} \BuiltInTok{sum}\NormalTok{(new\_bits[i] }\OperatorTok{\textless{}\textless{}}\NormalTok{ (n\_vars}\OperatorTok{{-}}\DecValTok{1}\OperatorTok{{-}}\NormalTok{i) }\ControlFlowTok{for}\NormalTok{ i }\KeywordTok{in} \BuiltInTok{range}\NormalTok{(n\_vars))}
\NormalTok{        result[new\_minterm] }\OperatorTok{=}\NormalTok{ (f\_bits }\OperatorTok{\textgreater{}\textgreater{}}\NormalTok{ minterm) }\OperatorTok{\&} \DecValTok{1}
    \ControlFlowTok{return} \BuiltInTok{sum}\NormalTok{(result[i] }\OperatorTok{\textless{}\textless{}}\NormalTok{ i }\ControlFlowTok{for}\NormalTok{ i }\KeywordTok{in} \BuiltInTok{range}\NormalTok{(}\DecValTok{2}\OperatorTok{**}\NormalTok{n\_vars))}

\KeywordTok{def}\NormalTok{ burnside\_count(function\_set, s4\_elements):}
    \CommentTok{"""Burnside 보조정리로 궤도 수 계산"""}
\NormalTok{    total\_fixed }\OperatorTok{=} \DecValTok{0}
    \ControlFlowTok{for}\NormalTok{ perm }\KeywordTok{in}\NormalTok{ s4\_elements:}
        \ControlFlowTok{for}\NormalTok{ f }\KeywordTok{in}\NormalTok{ function\_set:}
            \ControlFlowTok{if}\NormalTok{ apply\_permutation(f, perm) }\OperatorTok{==}\NormalTok{ f:}
\NormalTok{                total\_fixed }\OperatorTok{+=} \DecValTok{1}
    \ControlFlowTok{return}\NormalTok{ total\_fixed }\OperatorTok{//} \BuiltInTok{len}\NormalTok{(s4\_elements)}
\end{Highlighting}
\end{Shaded}

\begin{center}\rule{0.5\linewidth}{0.5pt}\end{center}

\subsubsection{3.7 차원 확장 (n=5) 계산
결과}\label{uxcc28uxc6d0-uxd655uxc7a5-n5-uxacc4uxc0b0-uxacb0uxacfc}

본 연구는 4변수(n=4) Boolean 함수 공간에 대한 완전한 분류를 수행했으나,
후속 연구의 일환으로 5변수(n=5) 공간에 대한 초기 확장을 진행하여 다음과
같은 결과를 도출했다.

{\def\LTcaptype{none} % do not increment counter
\begin{longtable}[]{@{}llll@{}}
\toprule\noalign{}
공간 & 함수 수 & S₅-궤도 수 & 검증 방법 \\
\midrule\noalign{}
\endhead
\bottomrule\noalign{}
\endlastfoot
Space(XOR)\_n5 & 64 & \textbf{12} & Burnside 이론값(12) 일치 \\
Space(AND)\_n5 & 7,581 & \textbf{210} & 직접 열거 + Burnside 일치 \\
\end{longtable}
}

\textbf{n=4 vs n=5 비교 요약:} - \textbf{Space(XOR)}: n=4에서 10개
궤도였으나, n=5에서 12개 궤도로 완만하게 증가했다. -
\textbf{Space(AND)}: n=4에서 30개 궤도였으나, n=5에서 210개 궤도로 크게
증가했다(Dedekind 수의 급증 반영).

\subsection{4. 선행 연구와의
통합}\label{uxc120uxd589-uxc5f0uxad6cuxc640uxc758-uxd1b5uxd569}

\subsubsection{4.1 정제 관계
(Refinement)}\label{uxc815uxc81c-uxad00uxacc4-refinement}

\textbf{정리 4.1 (정제 관계).} Repo 3의 S₄ 동치류 분류는 Repo 6의
Space(G) 분류보다 세밀한(finer) 분류이다.

\[\text{Space}(G) = \bigsqcup_{i} O_i \quad (S_4\text{-궤도들의 분리 합집합})\]

\begin{itemize}
\tightlist
\item
  같은 S₄-궤도에 속하는 두 함수 → 반드시 같은 Space(G)에 속함
\item
  같은 Space(G)에 속하는 두 함수 → 반드시 같은 S₄-궤도에 속할 필요 없음
\end{itemize}

\textbf{예시:} Space(XOR)에서 O₃ = \{A,B,C,D\}와 O₅ = \{A⊕B, A⊕C,
\ldots\}는 모두 Space(XOR) 안에 있지만 서로 다른 S₄-궤도에 속한다.

\subsubsection{4.2 Repo 1과의 통합: 체커보드의 군론적
설명}\label{repo-1uxacfcuxc758-uxd1b5uxd569-uxccb4uxcee4uxbcf4uxb4dcuxc758-uxad70uxb860uxc801-uxc124uxba85}

Repo 1의 핵심 발견 --- ``체커보드 패턴(A⊕B⊕C⊕D)은 변수 배열에 무관하게
동일한 K-map 패턴을 유지한다'' --- 은 이제 다음과 같이 군론으로 완전히
설명된다.

\begin{quote}
\textbf{체커보드 함수 \(A \oplus B \oplus C \oplus D\)는 Space(XOR)의
S₄-불동점이다.} S₄의 모든 24개 순열에 의해 고정되므로, 어떤 변수 배열을
선택해도 동일한 카르노 맵 패턴(O₉)이 나타난다.
\end{quote}

이것은 Repo 1에서 경험적으로 관찰된 불변성의 수학적 증명이다.

\subsubsection{4.3 Repo 2와의 통합: 대칭 함수의 군론적
설명}\label{repo-2uxc640uxc758-uxd1b5uxd569-uxb300uxce6d-uxd568uxc218uxc758-uxad70uxb860uxc801-uxc124uxba85}

Repo 2의 핵심 발견 --- ``대칭 Boolean 함수들이 Hamming weight 레이어
구조를 형성한다'' --- 은 다음과 같이 군론으로 설명된다.

\begin{quote}
\textbf{대칭적 단조 함수(threshold function) 6개는 Space(AND)의
S₄-불동점이다.} Hamming weight 레이어 구조는 이 함수들이 변수 치환에
불변이라는 사실의 시각적 표현이다(§3.2.3 참조).
\end{quote}

이것은 Repo 2에서 경험적으로 관찰된 레이어 구조의 수학적 근거다.

\subsubsection{4.4 Repo 3와의 통합: 변수 재배열 동치류와 S₄
궤도}\label{repo-3uxc640uxc758-uxd1b5uxd569-uxbcc0uxc218-uxc7acuxbc30uxc5f4-uxb3d9uxce58uxb958uxc640-sux2084-uxada4uxb3c4}

Repo 3의 핵심 작업 --- ``변수 재배열(S₄ 치환) 하의 동치류를 체계적으로
분류'' --- 은 본 논문의 S₄ 궤도 분류와 직접 대응한다.

\begin{quote}
\textbf{Repo 3의 동치류 ≡ 본 논문의 S₄ 궤도.} 다만 본 논문은 분류 대상을
Space(G)로 제한하고, Burnside 보조정리에 의한 정확한 궤도 수를
도출한다는 점에서 Repo 3의 분류를 수학적으로 정교화한다.
\end{quote}

이 관계는 §4.1의 정제 관계(정리 4.1)에서 이미 형식화되었다.

\subsubsection{4.5 통합 정리}\label{uxd1b5uxd569-uxc815uxb9ac}

{\def\LTcaptype{none} % do not increment counter
\begin{longtable}[]{@{}
  >{\raggedright\arraybackslash}p{(\linewidth - 4\tabcolsep) * \real{0.2143}}
  >{\raggedright\arraybackslash}p{(\linewidth - 4\tabcolsep) * \real{0.3929}}
  >{\raggedright\arraybackslash}p{(\linewidth - 4\tabcolsep) * \real{0.3929}}@{}}
\toprule\noalign{}
\begin{minipage}[b]{\linewidth}\raggedright
Repo
\end{minipage} & \begin{minipage}[b]{\linewidth}\raggedright
경험적 발견
\end{minipage} & \begin{minipage}[b]{\linewidth}\raggedright
군론적 해석
\end{minipage} \\
\midrule\noalign{}
\endhead
\bottomrule\noalign{}
\endlastfoot
Repo 1 & 체커보드 패턴의 변수 불변성 & \(A \oplus B \oplus C \oplus D\)
= Space(XOR)의 S₄-불동점 (§3.1.4) \\
Repo 2 & Hamming weight 레이어 구조 & 대칭적 단조 함수 = Space(AND)의
S₄-불동점 (§3.2.3) \\
Repo 3 & 변수 재배열 동치류 분류 & S₄ 궤도 = Burnside 계산 (§3 전체) \\
\end{longtable}
}

\subsection{5. 구조 인식 이론(SRT)과의
연결}\label{uxad6cuxc870-uxc778uxc2dd-uxc774uxb860srtuxacfcuxc758-uxc5f0uxacb0}

\subsubsection{5.1 SRT 핵심 가설
요약}\label{srt-uxd575uxc2ec-uxac00uxc124-uxc694uxc57d}

SRT(Repo 4)는 ``인간이 어떻게 새로운 구조를 인식하고 연구 질문을
생성하는가''를 설명하는 메타 이론이다. 본 논문과 관련된 핵심 가설:

{\def\LTcaptype{none} % do not increment counter
\begin{longtable}[]{@{}
  >{\raggedright\arraybackslash}p{(\linewidth - 2\tabcolsep) * \real{0.5000}}
  >{\raggedright\arraybackslash}p{(\linewidth - 2\tabcolsep) * \real{0.5000}}@{}}
\toprule\noalign{}
\begin{minipage}[b]{\linewidth}\raggedright
가설
\end{minipage} & \begin{minipage}[b]{\linewidth}\raggedright
내용
\end{minipage} \\
\midrule\noalign{}
\endhead
\bottomrule\noalign{}
\endlastfoot
H3 (주의 필터) & 인간은 모든 패턴 중 특정 구조에 먼저 주의를 기울인다 \\
H4 (의미 필터) & 인간은 ``연구할 가치 있는'' 구조를
압축성·예측력·일반화력으로 판단한다 \\
H5 (설명 예감) & 구조를 인식한 후 ``설명이 있을 것''이라는 예감이
발생한다 \\
H6 (설명적 유의미성) & 집합 전체의 유의미성 감지는 단일 함수의 유의미성
감지의 상위 수준이다 \\
H8 (개념-렌즈) & 획득한 개념이 이후 관찰의 렌즈가 된다 \\
H9 (개념 진화) & 개념은 연구를 통해 더 일반적인 형태로 진화한다 \\
H10 (생성적 개념) & 충분히 일반화된 개념은 새로운 연구 질문을 자율적으로
생성한다 \\
\end{longtable}
}

\subsubsection{5.2 S₄-불동점과 H3: 주의 필터의 수학적
설명}\label{sux2084-uxbd88uxb3d9uxc810uxacfc-h3-uxc8fcuxc758-uxd544uxd130uxc758-uxc218uxd559uxc801-uxc124uxba85}

\textbf{H3의 예측:} 인간은 65,536개의 K-map 패턴 중 특정 패턴에 먼저
주의를 기울인다.

\textbf{수학적 설명:} S₄-불동점 함수는 24가지 변수 배열
\textbf{모두}에서 동일한 패턴을 보인다. 반면 궤도 크기가 큰 함수는
대부분의 배열에서 다른 패턴을 보인다. 따라서, 임의의 변수 배열에서
카르노 맵을 관찰했을 때:

\begin{itemize}
\tightlist
\item
  \textbf{S₄-불동점 (궤도 크기 1)}: 어떤 배열에서도 동일한 패턴 →
  통계적으로 ``일관되게 눈에 띄는'' 함수
\item
  \textbf{궤도 크기 큰 함수 (크기 12, 24)}: 배열에 따라 다른 패턴 →
  상대적으로 덜 주목됨
\end{itemize}

\[\text{H3 예측 주의 순서: 아핀(체커보드) > 대칭(Maj) > 2-AND/2-OR > 임의 함수}\]

이는 궤도 크기(= 대칭 정도의 역수)가 낮을수록 더 많은 주의를 끈다는
의미이다.

\subsubsection{5.3 궤도 크기 스펙트럼과 H4: 의미 필터의
수치화}\label{uxada4uxb3c4-uxd06cuxae30-uxc2a4uxd399uxd2b8uxb7fcuxacfc-h4-uxc758uxbbf8-uxd544uxd130uxc758-uxc218uxce58uxd654}

H4는 ``연구할 가치 있는 구조를 판단하는 기준''으로
압축성(compressibility)을 제안한다. 궤도 크기는 이를 수치화하는
자연스러운 지표이다.

{\def\LTcaptype{none} % do not increment counter
\begin{longtable}[]{@{}
  >{\raggedright\arraybackslash}p{(\linewidth - 4\tabcolsep) * \real{0.3462}}
  >{\raggedright\arraybackslash}p{(\linewidth - 4\tabcolsep) * \real{0.1923}}
  >{\raggedright\arraybackslash}p{(\linewidth - 4\tabcolsep) * \real{0.4615}}@{}}
\toprule\noalign{}
\begin{minipage}[b]{\linewidth}\raggedright
궤도 크기
\end{minipage} & \begin{minipage}[b]{\linewidth}\raggedright
의미
\end{minipage} & \begin{minipage}[b]{\linewidth}\raggedright
Space 내 분포
\end{minipage} \\
\midrule\noalign{}
\endhead
\bottomrule\noalign{}
\endlastfoot
1 & S₄-불동점: 가장 대칭적, 가장 압축적 & XOR: 4개, AND: 6개, NOT:
\textbf{0개} \\
3\textasciitilde4 & 높은 대칭성 & AND: 11개 궤도 \\
6 & 중간 대칭성 & AND: 6개 궤도 \\
12\textasciitilde24 & 낮은 대칭성, 덜 압축적 & AND: 7개 궤도 \\
\end{longtable}
}

\textbf{Space(NOT)의 함의:} Space(NOT)에 S₄-불동점이 없다는 사실은
``자기쌍대이면서 완전히 대칭적인 함수''가 존재하지 않음을 의미한다. 이는
H4 의미 필터 관점에서, 자기쌍대 공간은 다른 공간(아핀, 단조)보다
``완벽히 규칙적인'' 구조를 포함하지 않는다는 흥미로운 함의를 가진다.

\subsubsection{5.4 H8-H9-H10: 개념 진화의
실사례}\label{h8-h9-h10-uxac1cuxb150-uxc9c4uxd654uxc758-uxc2e4uxc0acuxb840}

이 연구 프로그램 전체가 H8→H9→H10의 순차적 작동을 실사례로 보여준다.

\begin{verbatim}
체커보드 발견 [Repo 1]
  ↓ H3 작동 — "이 패턴이 눈에 띈다"
"불변성" 개념 형성
  ↓ H8 작동 — "불변성" 개념이 이후 관찰의 렌즈로 기능
층 구조·동치류 연구 [Repo 2, 3]
  ↓ H9 작동 — "불변성" → "S₄ 대칭" → "Space(G) 궤도"로 개념 진화
구조 인식 이론 [Repo 4]
  ↓ H10 작동 — SRT 자체가 새 연구(Repo 5, 6)를 생성
Boolean 함수 공간 이론 [Repo 6, 현재]
  ↓ 원점 복귀 — Repo 1의 체커보드가 이제 군론으로 완전히 설명됨
       ↓ 동시에 H10 재귀 작동 — Repo 6이 또 다른 새 연구를 생성
\end{verbatim}

특히 Repo 1의 체커보드 발견과 Repo 6의 S₄-불동점 분류가 \textbf{정확히
같은 함수를 지목}한다는 사실은 H5(설명 예감)의 소급적 증거이다: 연구자가
처음 ``왜?''라고 물었던 바로 그 함수가, 수학적으로도 ``가장 설명할 가치
있는'' 함수(S₄-불동점)임이 드러났다.

\begin{center}\rule{0.5\linewidth}{0.5pt}\end{center}

\subsection{6. 연구 기원
타임라인}\label{uxc5f0uxad6c-uxae30uxc6d0-uxd0c0uxc784uxb77cuxc778}

\emph{이 장은 연구의 인지 발달 과정을 SRT 프레임으로 정리한다. 개인 경험
섹션은 연구자가 직접 보완해야 한다.}

\subsubsection{Phase 0 --- 패턴 인식의
시작}\label{phase-0-uxd328uxd134-uxc778uxc2dduxc758-uxc2dcuxc791}

\begin{quote}
\textbf{{[}연구자 보완 필요{]}}\\
\emph{서술 가이드: 어린 시절(초등학생 시절 등) 수학적 패턴에서 처음 ``왜
이렇게 생겼을까''라는 질문이 발생한 구체적 경험을 서술해 주세요. 예:
구구단 표의 대칭, 바둑판 무늬, 타일 배열, 달력 숫자 패턴 등}
\end{quote}

\textbf{SRT 프레임:} 이 시기의 경험은 H3(주의 필터)의 최초 작동
사례이며, H1(구조 발견자 가설)의 원형이다.

\begin{center}\rule{0.5\linewidth}{0.5pt}\end{center}

\subsubsection{Phase 1 --- Karnaugh Map 체커보드 발견 (Repo
1)}\label{phase-1-karnaugh-map-uxccb4uxcee4uxbcf4uxb4dc-uxbc1cuxacac-repo-1}

\begin{quote}
\textbf{{[}연구자 보완 필요{]}}\\
\emph{서술 가이드: 카르노 맵을 처음 접한 상황(수업/과제/독서 중?),
체커보드 패턴에서 처음 ``이 패턴은 특별하다''는 인상을 받은 계기, 변수
배열을 바꿔도 패턴이 유지됨을 발견한 순간}
\end{quote}

\textbf{핵심 발견 (문서화됨):} 4변수 카르노 맵에서 XOR
함수(\(A \oplus B \oplus C \oplus D\))가 만드는 체커보드 패턴이 변수
배열 방식에 무관하게 유지된다. 현재(Repo 6)의 군론적 설명: 체커보드는
Space(XOR)의 S₄-불동점.

\textbf{SRT 프레임:} H3 최초 작동 (주의 포착) → H5 발동 (설명 예감).

\begin{center}\rule{0.5\linewidth}{0.5pt}\end{center}

\subsubsection{Phase 2 --- 대칭 Boolean 함수의 시각 패턴 (Repo
2)}\label{phase-2-uxb300uxce6d-boolean-uxd568uxc218uxc758-uxc2dcuxac01-uxd328uxd134-repo-2}

\begin{quote}
\textbf{{[}연구자 보완 필요{]}}\\
\emph{서술 가이드: Phase 1에서 Phase 2로 넘어간 개념적 연결이 어떻게
의식적으로 형성되었는가? ``층 구조'' 개념이 처음 등장한 순간}
\end{quote}

\textbf{핵심 발견 (문서화됨):} 대칭 Boolean 함수들이 K-map에서 Hamming
Weight 레이어 구조를 형성한다. 현재(Repo 6)의 군론적 설명: 대칭 함수 =
Space(AND)의 S₄-불동점(단조인 경우).

\textbf{SRT 프레임:} H4 작동 (의미 판단) → H8 전조 (``층 구조''가 렌즈로
기능하기 시작).

\begin{center}\rule{0.5\linewidth}{0.5pt}\end{center}

\subsubsection{Phase 3 --- 변수 재배열 불변성 이론화 (Repo
3)}\label{phase-3-uxbcc0uxc218-uxc7acuxbc30uxc5f4-uxbd88uxbcc0uxc131-uxc774uxb860uxd654-repo-3}

\begin{quote}
\textbf{{[}연구자 보완 필요{]}}\\
\emph{서술 가이드: 구체적인 수학적 돌파구(R-3, A-7 등)가 언제, 어떤
맥락에서 도출되었는가? ``동치(Equivalence)''라는 개념이 처음 명시적으로
등장한 순간}
\end{quote}

\textbf{핵심 발견 (문서화됨):} 변수 재배열(S₄) 하에서 65,536개 함수를
동치류로 분류. 현재(Repo 6)와의 관계: 정제 관계(Theorem 4.1) --- Repo
3의 분류는 Repo 6 분류의 세밀한 버전.

\textbf{SRT 프레임:} H5·H6 작동 → H8 본격화 (``불변성'' 개념이 연구
전체를 보는 렌즈로 기능).

\begin{center}\rule{0.5\linewidth}{0.5pt}\end{center}

\subsubsection{Phase 4 --- 구조 인식 이론 탄생 (Repo
4)}\label{phase-4-uxad6cuxc870-uxc778uxc2dd-uxc774uxb860-uxd0c4uxc0dd-repo-4}

\begin{quote}
\textbf{{[}연구자 보완 필요{]}}\\
\emph{서술 가이드: Boolean 함수 연구에서 인지 이론으로의 도약이 일어난
구체적 계기. ``왜 이것이 연구할 가치가 있어 보이는가?''라는 메타 질문이
처음 등장한 순간}
\end{quote}

\textbf{핵심 전환 (문서화됨):} Boolean 함수의 특수한 성질 연구에서,
인간이 어떻게 연구 가치 있는 구조를 발견하는가라는 일반 이론으로의 도약.

\textbf{SRT 프레임:} H9 작동 --- ``불변성'' 개념이 ``구조 인식''으로
진화.

\begin{center}\rule{0.5\linewidth}{0.5pt}\end{center}

\subsubsection{Phase 5 --- 인간-AI 협업의 방법론적 성찰 (Repo
5)}\label{phase-5-uxc778uxac04-ai-uxd611uxc5c5uxc758-uxbc29uxbc95uxb860uxc801-uxc131uxcc30-repo-5}

\textbf{핵심 전환 (문서화됨):} 연구 방법(AI와 함께 연구하는 방식) 자체가
연구 대상이 됨.

{\def\LTcaptype{none} % do not increment counter
\begin{longtable}[]{@{}ll@{}}
\toprule\noalign{}
인간의 역할 & AI의 역할 \\
\midrule\noalign{}
\endhead
\bottomrule\noalign{}
\endlastfoot
이상한 현상 발견 & 개념 연결 \\
이상 징후 감지 & 프레임워크 제안 \\
질문 생성 & 설명 공간 확장 \\
\end{longtable}
}

\textbf{SRT 프레임:} H10 작동 --- SRT 자체가 새 연구(Repo 5, 6)를
생성하는 개념으로 기능.

\begin{center}\rule{0.5\linewidth}{0.5pt}\end{center}

\subsubsection{Branch 8 --- Boolean 함수 공간 이론 (Repo 6,
현재)}\label{branch-8-boolean-uxd568uxc218-uxacf5uxac04-uxc774uxb860-repo-6-uxd604uxc7ac}

\textbf{핵심 동기 (문서화됨):} Phase 1의 체커보드 발견이 이제 군론(S₄
궤도 분류)으로 완전히 설명됨. 최초 발견에서 수학적 통합까지의 원점
복귀이자, 새로운 수학적 확장의 출발점.

\textbf{SRT 프레임:} - S₄-불동점(완전 대칭 함수들) = Phase 1에서 처음
``눈에 띄었던'' 바로 그 함수들 → H3·H5 소급 검증 - H10의 재귀적 작동 ---
Repo 6이 또 다른 새 연구(차원 확장, 인지 실험 등)를 생성

\begin{center}\rule{0.5\linewidth}{0.5pt}\end{center}

\subsection{7. 결론}\label{uxacb0uxb860}

\subsubsection{7.1 수학적 기여
요약}\label{uxc218uxd559uxc801-uxae30uxc5ec-uxc694uxc57d}

본 연구는 4변수 Boolean 함수 공간의 S₄ 군론적 완전 분류를 수행하였다. 세
주요 결과:

\begin{enumerate}
\def\labelenumi{\arabic{enumi}.}
\item
  \textbf{Space(XOR) 10궤도}: Burnside 보조정리와 Python 컴퓨터 열거로
  이중 검증. S₄-불동점 4개이며, 이 중
  체커보드(\(A \oplus B \oplus C \oplus D\))가 Repo 1의 연구 출발점과
  정확히 일치함을 확인.
\item
  \textbf{Space(AND) 30궤도}: 168개 단조 함수의 완전 분류. S₄-불동점
  6개(임계값 함수)가 Repo 2의 대상과 정확히 일치. 궤도 크기 분포(1×6,
  3×2, 4×9, 6×6, 12×7)의 완전 목록 수립.
\item
  \textbf{Space(NOT) 32궤도 및 불동점 부재 증명(정리 3.1)}:
  ``자기쌍대이면서 완전히 대칭적인'' 함수는 4변수에서 존재하지 않는다는
  음성 결과. HW=2 레이어의 보수 쌍 충돌 메커니즘으로 증명.
\end{enumerate}

\subsubsection{7.2 이론적 기여
요약}\label{uxc774uxb860uxc801-uxae30uxc5ec-uxc694uxc57d}

수학적 분류가 선행 연구 3편(Repo 1--3)의 경험적 발견을 단일 군론
프레임워크로 통합한다(정제 관계, 정리 4.1). 또한 SRT의 핵심 가설들에
구체적인 수학적 근거를 제공한다:

\begin{itemize}
\tightlist
\item
  \textbf{H3 설명}: S₄-불동점 함수들이 통계적으로 가장 많은 배열에서
  동일한 패턴을 보이므로, ``먼저 눈에 띄는'' 것은 수학적으로 자연스럽다.
\item
  \textbf{H4 수치화}: 궤도 크기가 낮을수록 압축성이 높으며, 이것이
  ``연구 가치'' 판단의 수학적 기반이다.
\item
  \textbf{H8·H9 실사례}: 이 연구 프로그램 전체의 개념 발전 경로(체커보드
  → 불변성 → S₄ 대칭 → 구조 인식)가 H8→H9→H10의 순차 작동을 보여준다.
\end{itemize}

\subsubsection{7.3 미해결 문제와 향후
연구}\label{uxbbf8uxd574uxacb0-uxbb38uxc81cuxc640-uxd5a5uxd6c4-uxc5f0uxad6c}

\paragraph{7.3.1 미해결 수학적
문제}\label{uxbbf8uxd574uxacb0-uxc218uxd559uxc801-uxbb38uxc81c}

\textbf{문제 7.1 (S₄-부분공간의 완전 분류):} \(\mathcal{F}\) 안의 모든
S₄-부분공간을 완전히 분류하라.

현재 알려진 S₄-부분공간: - 자명한 것: \(\emptyset\), \(\{0\}\),
\(\{1\}\), \(\{0, 1\}\), \(\mathcal{F}\) - L (XOR 공간), M (AND 공간), L
∩ M - L의 궤도들: \{0\}, \{1\}, \{x₁,x₂,x₃,x₄\} 등

\textbf{문제 7.2 (구조 복잡성과 인지 용이성의 관계):} 궤도 크기 s와
인간의 구조 인식 시간 T 사이의 정량적 관계를 밝혀라.

\textbf{가설:} \(T \propto \log(s)\) (대수적 관계)

\paragraph{7.3.2 실험적 검증 계획 및 향후 연구
로드맵}\label{uxc2e4uxd5d8uxc801-uxac80uxc99d-uxacc4uxd68d-uxbc0f-uxd5a5uxd6c4-uxc5f0uxad6c-uxb85cuxb4dcuxb9f5}

\begin{quote}
향후 연구의 전체 로드맵과 상세 설계는 \href{TODO.md}{\texttt{TODO.md}}의
``미래 연구(FutureWorks)'' 절에 통합되어 있다.
\end{quote}

현재 Repo 6의 Phase 3 (실험 단계): - 자극 설계: 10개 S₄-궤도 대표 원소
(크기 1, 4, 6, 12 각 2-3개) - 측정: 패턴 인식 시간, 정확도, 주관적
복잡성 평가 - 피험자: N=30 (IRB 승인 후)

{\def\LTcaptype{none} % do not increment counter
\begin{longtable}[]{@{}
  >{\raggedright\arraybackslash}p{(\linewidth - 4\tabcolsep) * \real{0.3750}}
  >{\raggedright\arraybackslash}p{(\linewidth - 4\tabcolsep) * \real{0.3750}}
  >{\raggedright\arraybackslash}p{(\linewidth - 4\tabcolsep) * \real{0.2500}}@{}}
\toprule\noalign{}
\begin{minipage}[b]{\linewidth}\raggedright
우선순위
\end{minipage} & \begin{minipage}[b]{\linewidth}\raggedright
연구 방향
\end{minipage} & \begin{minipage}[b]{\linewidth}\raggedright
상태
\end{minipage} \\
\midrule\noalign{}
\endhead
\bottomrule\noalign{}
\endlastfoot
1 & \textbf{SRT 인지 실험 (Phase 3)} --- H3·H4·H6 수학적 예측의 실증
검증 & ⏸️ 후속 연구 이관 \\
2 & \textbf{인간-AI 협업 모델·교육 프레임워크 (Phase 4)} --- 탐구 루프와
Level 1\textasciitilde3 시나리오 & ⏸️ 후속 연구 이관 \\
3 & \textbf{차원 확장 (n=5)} --- S₅-궤도 완전 계산 완료 & ✅ 완료
(Space(XOR)=12, Space(AND)=210) \\
4 & \textbf{다치 논리 확장} --- k치 논리 Space(G) 일반화 & 장기 단계 \\
5 & \textbf{연구 기원 타임라인 완성} --- paper.md §6 개인 서술 보완 &
연구자 직접 작성 필요 \\
\end{longtable}
}

\begin{center}\rule{0.5\linewidth}{0.5pt}\end{center}

\subsection{참고 저장소 및 관련
문서}\label{uxcc38uxace0-uxc800uxc7a5uxc18c-uxbc0f-uxad00uxb828-uxbb38uxc11c}

{\def\LTcaptype{none} % do not increment counter
\begin{longtable}[]{@{}
  >{\raggedright\arraybackslash}p{(\linewidth - 4\tabcolsep) * \real{0.2692}}
  >{\raggedright\arraybackslash}p{(\linewidth - 4\tabcolsep) * \real{0.1923}}
  >{\raggedright\arraybackslash}p{(\linewidth - 4\tabcolsep) * \real{0.5385}}@{}}
\toprule\noalign{}
\begin{minipage}[b]{\linewidth}\raggedright
저장소
\end{minipage} & \begin{minipage}[b]{\linewidth}\raggedright
내용
\end{minipage} & \begin{minipage}[b]{\linewidth}\raggedright
이 논문과의 연결
\end{minipage} \\
\midrule\noalign{}
\endhead
\bottomrule\noalign{}
\endlastfoot
\href{https://inha20.github.io/1KMapStructureInvariance/}{1KMapStructureInvariance}
& K-map 패턴 불변성 발견 & 체커보드 = Space(XOR) O₉ (§3.1.4, §4.2) \\
\href{https://inha20.github.io/2SymmetricBooleanFunctionMinorThesis/}{2SymmetricBooleanFunctionMinorThesis}
& 대칭 함수 K-map 패턴 & 임계값 함수 = Space(AND) S₄-불동점 (§3.2.3,
§4.3) \\
\href{https://inha20.github.io/3VariableRearrangementInvarianceMinorThesis/}{3VariableRearrangementInvarianceMinorThesis}
& 변수 재배열 동치류 분류 & 정제 관계 (§4.1, §4.4) \\
\href{https://inha20.github.io/4StructureRecognitionTheory/}{4StructureRecognitionTheory}
& SRT: H1--H10, Def. 3.1--3.3 & SRT와의 연결 (§5) \\
\href{https://inha20.github.io/5HumanAIResearchCollaboration/}{5HumanAIResearchCollaboration}
& 인간-AI 협업 방법론 & 연구 기원 타임라인 §Phase 5 \\
\texttt{theory/S4GroupAnalysis.md} & 핵심 수학 문서
(OC-1\textasciitilde OC-5 완전 분류) & §3 전체 \\
\texttt{scripts/s4\_orbit\_calculator.py} & 궤도 계산 Python 스크립트
(n=4 전체 + n=5 확장) & 컴퓨터 검증 \\
\texttt{TODO.md} & 마스터 로드맵 + 미래 연구(FutureWorks,
FR-1\textasciitilde FR-3) & §7.3 연결 \\
\end{longtable}
}

\begin{center}\rule{0.5\linewidth}{0.5pt}\end{center}

\emph{본 논문은 \texttt{6BooleanFunctionSpaceTheory} 저장소의 Phase 1·2
완료 시점을 기준으로 작성된 통합 본문 문서입니다.}\\
\emph{Phase 3(인지 실험)은 후속 연구로 이관되었습니다. §6의
\texttt{{[}연구자\ 보완\ 필요{]}} 섹션은 연구자가 직접 채워야 합니다.}\\
\emph{관련 이론: \texttt{StructureRecognitionTheory\_Unified.md}
H8·H9·H10, Def. 3.1--3.3}

\end{document}
